\section{Evaluation}
\label{sec:evaluation}

We evaluate \axvisor{} along four research questions that together test whether
substrate-OS portability is a real, useful property rather than a veneer over
per-OS reimplementation.

\begin{description}
  \item[RQ1 (Portability cost).] How much substrate-specific code is required to
  host the \axvisor{} core on each substrate, and how much of the core is
  shared?
  \item[RQ2 (Functionality).] Does the contract suffice to run a real,
  unmodified guest workload on each substrate?
  \item[RQ3 (Guest-side neutrality).] Does the KVM-compatible ABI let the same
  guest and toolchain run unchanged across substrates?
  \item[RQ4 (Overhead).] What virtualization overhead does the portable design
  impose relative to a non-portable baseline, and how does overhead vary across
  substrates?
\end{description}

\subsection{Experimental Setup}

We deploy \axvisor{} on three substrates: \arceos{}, \asterinas{}, and Linux,
each running on \fillme{target board / QEMU version and ISA, e.g., QEMU virt
ARMv8-A or x86-64}. The guest is an unmodified bare-metal Linux kernel
(\fillme{version}) booted via the KVM-compatible ABI exposed by the core. The
baseline for overhead comparisons is the original \axvisor{}-on-\arceos{}
configuration prior to refactoring, which represents the non-portable,
substrate-fused design; where useful we also compare against native KVM on Linux
as a maturity reference. All measurements are taken on \fillme{hardware, CPU,
memory, core count}. Each data point is the median of \fillme{N} runs; we report
\fillme{metric for dispersion}. Numerical results marked with \fillme{} are to
be filled from the authors' measurement campaign; the methodology and structure
below are final.

\subsection{RQ1: Portability Cost}

Portability cost has two facets: the size of each adaptation layer and the
fraction of the system that is shared. We measure both. We count lines of code
in the adaptation layer for each substrate and in the portable core, and we
compute the core-reuse ratio, defined as the fraction of total hypervisor code
(core + all adaptation layers) that resides in the shared core.

Table~\ref{tab:portability-cost} reports the breakdown. The expectation from the
design is that the core dominates and each adaptation layer is small relative to
it, because adaptation layers contain only contract implementations rather than
virtualization logic. The \arceos{} layer should be the smallest, since
\arceos{} is already hypervisor-friendly; the Linux layer the largest, because
bridging Linux's abstractions requires the most reconciliation; and \asterinas{}
in between.

\begin{table}[t]
  \centering
  \caption{Portability cost: code distribution across the core and adaptation
  layers. \fillme{Fill with measured values}.}
  \label{tab:portability-cost}
  \begin{tabular}{@{}lccp{1.0cm}l@{}}
    \toprule
    \textbf{Component} & \textbf{LoC} & \textbf{\% total} &
      \textbf{\# ops} & \textbf{Notes} \\
    \midrule
    Portable core   & \fillme{} & \fillme{} & ---        & shared by all \\
    \arceos{} layer & \fillme{} & \fillme{} & \fillme{}  & baseline \\
    \asterinas{} layer & \fillme{} & \fillme{} & \fillme{} & framekernel \\
    Linux layer     & \fillme{} & \fillme{} & \fillme{}  & monolithic \\
    \midrule
    \multicolumn{2}{@{}l}{Core reuse ratio} &
      \multicolumn{3}{l}{\fillme{X\%}} \\
    \bottomrule
  \end{tabular}
\end{table}

Beyond raw size, we analyze \emph{where} adaptation effort concentrates. We
classify each contract operation by the kind of work its adaptation required:
\emph{direct mapping} (the substrate API already matches the contract),
\emph{bridging} (the substrate provides the capability but behind a different
abstraction), and \emph{recovery} (the substrate hides the needed resource and
the layer must reconstruct it). We expect memory and interrupt operations to
dominate bridging and recovery, consistent with the semantic mismatches
described in \S\ref{sec:adaptation-layers}. This distribution is itself a
finding: it shows that portability cost is not uniform across the contract but
clusters where OS designs diverge most, validating the contract's structure as a
faithful map of where portability is contested.

\subsection{RQ2: Functionality}

To test whether the contract is sufficient, we boot an unmodified bare-metal
Linux guest on each substrate and exercise the virtualization paths it implies:
VM creation, vCPU run/pause, stage-2 memory mapping for the guest's physical
memory, MMIO trap handling for emulated devices, interrupt injection for the
guest's timer and console, and hypercall handling.
Table~\ref{tab:functionality} records the functionality matrix.

\begin{table}[t]
  \centering
  \caption{Functionality matrix: guest capabilities supported per substrate via
  the contract. \fillme{Confirm per substrate}.}
  \label{tab:functionality}
  \begin{tabular}{@{}lccc@{}}
    \toprule
    \textbf{Capability} & \textbf{\arceos{}} & \textbf{\asterinas{}} &
      \textbf{Linux} \\
    \midrule
    Guest boot (bare-metal Linux) & \fillme{} & \fillme{} & \fillme{} \\
    vCPU create/run               & \fillme{} & \fillme{} & \fillme{} \\
    Stage-2 memory mapping        & \fillme{} & \fillme{} & \fillme{} \\
    MMIO trap \& emulate          & \fillme{} & \fillme{} & \fillme{} \\
    Interrupt injection           & \fillme{} & \fillme{} & \fillme{} \\
    Timer virtualization          & \fillme{} & \fillme{} & \fillme{} \\
    Virtio console                & \fillme{} & \fillme{} & \fillme{} \\
    \bottomrule
  \end{tabular}
\end{table}

The decisive result is that the same guest image boots on both the \asterinas{}
and Linux substrates, which differ substantially in kernel architecture. This
demonstrates that the contract captures the capabilities required for a real
guest, not merely for a trivial loop. Where a capability is absent on a
substrate, the cause is a missing optional capability (e.g., IOMMU passthrough)
rather than a contract insufficiency, and the core degrades rather than
failing---consistent with the capability-based design of \S\ref{sec:contract}.

\subsection{RQ3: Guest-Side Neutrality}

RQ3 asks whether the KVM-compatible ABI achieves guest-side decoupling. We test
this by using the same guest image, the same guest kernel configuration, and the
same VMM invocation sequence on the \asterinas{} and Linux substrates. If the
ABI shim is correct and complete, no guest-side change should be needed when
switching substrates.

We verify this in two ways. First, we confirm that the guest boots and reaches
userspace on both substrates with identical images. Second, we exercise a
representative set of KVM ioctls and check that their behavior matches the
expected semantics. The result, summarized in Table~\ref{tab:abi-coverage},
shows that the guest and toolchain are invariant across substrates: the
substrate choice is invisible above the core. This is the practical payoff of
double decoupling---the guest $\times$ substrate matrix is covered without
per-cell work.

\begin{table}[t]
  \centering
  \caption{KVM-compatible ABI coverage and guest-side invariance across
  substrates. \fillme{Fill with measured coverage}.}
  \label{tab:abi-coverage}
  \begin{tabular}{@{}lcc@{}}
    \toprule
    \textbf{Property} & \textbf{\asterinas{}} & \textbf{Linux} \\
    \midrule
    Identical guest image boots  & \fillme{yes}  & \fillme{yes} \\
    Identical VMM/toolchain used & \fillme{yes}  & \fillme{yes} \\
    KVM ioctl coverage (\%)      & \fillme{}     & \fillme{} \\
    Guest-side changes required  & \fillme{none} & \fillme{none} \\
    \bottomrule
  \end{tabular}
\end{table}

\subsection{RQ4: Virtualization Overhead}

Finally, we measure the cost that the portable design imposes. We compare three
configurations: the original non-portable \axvisor{}-on-\arceos{} (baseline),
the refactored portable \axvisor{} on each substrate, and native KVM on Linux as
a maturity reference (not as a strict fairness baseline, since native KVM is far
more mature). We report four classes of metric.

\paragraph{VM lifecycle}
Guest boot time and VM teardown time, measured from \code{KVM\_CREATE\_VM} to
guest userspace ready. This exercises memory mapping and vCPU creation paths.

\paragraph{Trap handling}
VM-exit/entry round-trip latency for representative exit reasons (MMIO,
hypercall, timer). This is the most direct measure of core path overhead, since
it is dominated by the \code{vcpu\_run} loop and trap dispatch that live in the
shared core.

\paragraph{Memory operations}
Latency of stage-2 map/unmap and TLB maintenance, which exercise the memory
group of the contract and therefore reflect adaptation-layer efficiency.

\paragraph{I/O path}
Throughput and latency of the virtio console / virtio block path if available,
exercising MMIO emulation and interrupt injection together.

\begin{table}[t]
  \centering
  \caption{Virtualization overhead across substrates. Lower is better except
  where noted. \fillme{Fill with measured values}.}
  \label{tab:overhead}
  \begin{tabular}{@{}lcccc@{}}
    \toprule
    \textbf{Metric} & \textbf{Baseline} & \textbf{\arceos{}} &
      \textbf{\asterinas{}} & \textbf{Linux} \\
    \midrule
    Guest boot (ms)     & \fillme{} & \fillme{} & \fillme{} & \fillme{} \\
    VM-exit RT ($\mu$s) & \fillme{} & \fillme{} & \fillme{} & \fillme{} \\
    Stage-2 map ($\mu$s)& \fillme{} & \fillme{} & \fillme{} & \fillme{} \\
    TLB flush ($\mu$s)  & \fillme{} & \fillme{} & \fillme{} & \fillme{} \\
    I/O latency ($\mu$s)& \fillme{} & \fillme{} & \fillme{} & \fillme{} \\
    \bottomrule
  \end{tabular}
\end{table}

The expected outcome, supported by the architecture, is that the portable
design's overhead is concentrated in adaptation-layer bridging for memory and
interrupt operations, while trap-handling latency---which stays in the shared
core---should be close to the baseline on every substrate. Variation across
substrates reflects the cost of reconciling each OS's abstractions, not the cost
of portability in the abstract: the same core runs on all three, so
cross-substrate differences isolate substrate-induced overhead. If a substrate's
overhead is higher, the cause is traceable to a specific contract operation's
adaptation, which is itself actionable feedback for that substrate rather than a
deficiency of the core.

\subsection{Threats to Validity}

Three threats bear noting. First, three substrates cannot exhaust the OS design
space; \arceos{}, \asterinas{}, and Linux span unikernel, framekernel, and
monolithic models, but microkernel and capability-based substrates remain future
work. Second, the guest workload---a bare-metal Linux boot---is representative
but not comprehensive; richer guests and I/O-intensive workloads would stress
the contract further. Third, the overhead comparison against native KVM is
asymmetric in maturity; we therefore treat KVM as a reference point, not a
fairness-controlled baseline, and focus the comparison on the
portable-vs.-non-portable \axvisor{} contrast that isolates the design's own
cost.
